% Generated by tools/edn_latex.py from reports/edn.md.
% Do not edit: change reports/edn.md and run `make edn`.
\ednentry{
  id = {EDN-12},
  title = {Retraining and promotion are human-triggered, not automated},
  date = {2026-09-22},
  milestone = {M6: Monitoring},
  topic = {ML System Design, Model Performance},
  participants = {@lukas2510},
  decision = {FR-15 formalises the retrain-after-drift loop that EDN-03 was designed to demonstrate: when FR-14's drift job flags input drift on the \texttt{ES} replay (per NFR-11), a team member retrains with \texttt{ES} included and opens the promotion PR by hand. No automation watches the drift job and triggers this on its own. The candidate still must pass NFR-01's gate before the PR is merged; merging is the promotion (EDN-08), and rollback is redeploying the previous image tag (EDN-08).},
  rationale = {\emph{Alternatives considered:}
\begin{itemize}[leftmargin=*, topsep=2pt]
\item \textbf{Option A: fully automated retraining and promotion.} A watcher on the drift job's output triggers \texttt{dvc repro} and opens the promotion PR automatically when drift is confirmed.
\newline Pros: reads as more ``extended/innovative'' for the M6 System Design criterion; no manual step to forget.
\newline Cons: needs new infrastructure (something connecting Alibi Detect's output to a CI trigger) built and debugged in the last two weeks before the Dec 9 presentation, exactly when there is the least slack to recover if it breaks; higher risk for a 5-person team with no prior automation of this kind in the project.
\item \textbf{Option B (chosen): human-triggered retraining and promotion.} The drift job surfaces the signal (Grafana); a team member runs the retrain and opens the PR once NFR-01's gate passes.
\newline Pros: still fully closes the loop EDN-03 was built to demonstrate, and still satisfies the M6 rubric's ``enabling confident investigation and response'', a human response is a legitimate response, not a lesser one; far less new infrastructure, and safer to demo live; matches the team's size and the amount of time left before M6.
\newline Cons: less visually impressive than full automation; a person has to actually notice the signal and act on it.
\end{itemize}
\par
\emph{Why:} The \texttt{ES} holdout exists specifically to demonstrate this loop closing (EDN-03), so leaving it as an unformalised sentence in the brief risked it quietly not getting built once M6 crunch hits. Given the timeline, automating the trigger is a real risk for a low grading upside over the human-triggered version, which already gets full credit for ``closing the loop'' and ``responding to a signal'' without the extra failure surface this late in the semester.},
  ai = {Information seeking, Alternative generation, Alternative assessment, Recommendation},
  response = {Accepted},
  assessment = {While reviewing PR \#19, AI noted that the retrain-after-drift plan in the brief had never become a testable requirement, unlike everything else on the requirements page, and that EDN-08 had already defined the promotion and rollback mechanism without it being referenced by any FR. Asked whether this and a related NFR-11 tightening were still relevant and worth doing, AI confirmed both were unaddressed, proposed the NFR-11 fix directly (no real alternative to weigh), and for retraining, named the automation-level choice as a genuine decision, laid out both options with the project's timeline as the deciding factor, and recommended the human-triggered option. Lukas confirmed it.},
  aievidence = {Claude Code session on 2026-09-22 while reviewing PR \#19: AI's original review had flagged both the trivially-passable NFR-11 and the missing retraining requirement; Lukas asked whether they were still relevant and whether they were overkill, AI confirmed relevance for both and recommended the human-triggered option for retraining with reasoning tied to the M6 timeline; Lukas replied ``yes.''},
  evidence = {\href{https://github.com/mlops-2627q1-mds-upc/MLOPS_Recommenditos/blob/main/docs/docs/requirements.md}{requirements} FR-14, FR-15, NFR-01, NFR-11; \href{https://github.com/mlops-2627q1-mds-upc/MLOPS_Recommenditos/blob/main/docs/docs/specification.md}{specification} FR-14 and FR-15; \href{https://github.com/mlops-2627q1-mds-upc/MLOPS_Recommenditos/blob/main/docs/docs/project-brief.md}{project brief} section 3.2 (\texttt{ES} new-market drift scenario) and section 5 (target architecture); \hyperref[EDN-03]{EDN-03}, \hyperref[EDN-08]{EDN-08}; \href{https://github.com/mlops-2627q1-mds-upc/MLOPS_Recommenditos/pull/19}{PR \#19}.},
}
